% Part: quantified-modal-logic
% Chapter: syntax-and-semantics
% Section: languages

\documentclass[../../../include/open-logic-section]{subfiles}

\begin{document}

\olfileid{qml}{syn}{lan}

\olsection{Language}

\begin{explain}
The language of quantified modal logic combines those of
first-order logic (terms, predicates, quantification)
and modal logic (modal operators).
\end{explain}

Expressions of quantified modal logic are built up from a basic vocabulary
containing \emph{!!{variable}s}, \emph{!!{constant}s},
\emph{!!{predicate}s} and sometimes \emph{!!{function}s}.  From them,
together with logical connectives, quantifiers, modal operators, 
and punctuation symbols such as parentheses and commas, \emph{terms} and
\emph{!!{formula}s} are formed. 

\begin{enumerate}
\item Logical symbols
\begin{enumerate}
\item Logical connectives:
  \startycommalist
  \iftag{prvFalse}{\ycomma $\lfalse$ (the propositional constant for !!{falsity})}{}%
  \iftag{prvTrue}{\ycomma $\ltrue$ (the propositional constant for !!{truth})}{}%
  \iftag{prvNot}{\ycomma $\lnot$ (negation)}{}%
  \iftag{prvAnd}{\ycomma $\land$ (conjunction)}{}%
  \iftag{prvOr}{\ycomma $\lor$ (disjunction)}{}%
  \iftag{prvIf}{\ycomma $\lif$ (!!{conditional})}{}%
  \iftag{prvIff}{\ycomma $\liff$ (!!{biconditional})}{}%
  \iftag{prvAll}{\ycomma $\lforall$ (universal quantifier)}{}%
  \iftag{prvEx}{\ycomma $\lexists$ (existential quantifier)}{}%
  \iftag{prvBox}{\ycomma $\Box$ (necessity operator)}{}%
  \iftag{prvDiamond}{\ycomma $\Diamond$ (possibility operator)}{}.
\item The two-place !!{identity}~$\eq$.
\item A !!{denumerable}s set of !!{variable}s: $\Obj v_0$, $\Obj v_1$, $\Obj
  v_2$, \dots
\end{enumerate}
\item Non-logical symbols, making up the \emph{standard
  language} of quantified modal logic
\begin{enumerate}
\item A !!{denumerable}s set of $n$-place !!{predicate}s for each $n>0$: $\Obj
  A^n_0$, $\Obj A^n_1$, $\Obj A^n_2$, \dots
\item A !!{denumerable}s set of !!{constant}s: $\Obj c_0$, $\Obj c_1$, $\Obj
  c_2$, \dots.
\item A !!{denumerable}s set of $n$-place !!{function}s for each $n>0$:
  $\Obj f^n_0$, $\Obj f^n_1$, $\Obj f^n_2$, \dots
\end{enumerate}
\item Punctuation marks: (, ), and the comma.
\end{enumerate}

Most of our definitions and results will be formulated for the full
standard language.  However, depending on the
application, we may also restrict the language to only a few
!!{predicate}s, !!{constant}s, and !!{function}s.

\iftag{defNot,defOr,defAnd,defIf,defIff,defTrue,defFalse,defEx,defAll}{%
In addition to the primitive connectives and
\iftag{notprvEx,notprvAll}{quantifier}{quantifiers} introduced
above, we also use the following \emph{defined} symbols:
\startycommalist
  \iftag{defFalse}{\ycomma !!{falsity}~$\lfalse$}{}%
  \iftag{defTrue}{\ycomma !!{truth}~$\ltrue$}}{}%
  \iftag{defNot}{\ycomma $\lnot$ (negation)}{}%
  \iftag{defAnd}{\ycomma $\land$ (conjunction)}{}%
  \iftag{defOr}{\ycomma $\lor$ (disjunction)}{}%
  \iftag{defIf}{\ycomma $\lif$ (!!{conditional})}{}%
  \iftag{defIff}{\ycomma $\liff$ (!!{biconditional})}{}%
  \iftag{defAll}{\ycomma $\lforall$ (universal quantifier)}{}%
  \iftag{defEx}{\ycomma $\lexists$ (existential quantifier)}{}%
  \iftag{defBox}{\ycomma $\Box$ (necessity operator)}{}%
  \iftag{defDiamond}{\ycomma $\Diamond$ (possibility operator)}{}.

\begin{tagblock}{defNot,defOr,defAnd,defIf,defIff,defTrue,defFalse,defEx,defAll,defBox,defDiamond}
\begin{explain}
A defined symbol is not officially part of the language, but is
introduced as an informal abbreviation: it allows us to abbreviate
formulas which would, if we only used primitive symbols, get quite
long.  This is obviously an advantage.  The bigger advantage, however,
is that proofs become shorter.  If a symbol is primitive, it has to be
treated separately in proofs. The more primitive symbols, therefore,
the longer our proofs.
\end{explain}
\end{tagblock}


\end{document}
